\documentclass[a4paper, 12pt]{article}

\usepackage[utf8]{inputenc}
\usepackage{enumitem}
\usepackage{amsmath}
\usepackage{amssymb}

\title{E.B.1.2}
\author{Lorenzo Burello}
\date{\today}

\begin{document}
\maketitle

\section{Definizione del problema}
\subsection{Definizione FOL}
Si consideri $\mathcal{P}$:
\begin{itemize}[label=-]
    \item $Studente/1$
    \item $Corso/1$
    \item $HaStudiato/2$
    \item $Ansioso/1$
    \item $Supera/2$
\end{itemize}
E la formula $\phi$:
\begin{equation}
    \forall X \forall C \;\; Studente(X)\land Corso(C) \land (Ansioso(X) \lor \lnot HaStudiato(X,C)) \rightarrow \lnot Supera(X, C)
\end{equation}
\subsection{Inferenza A}
Si dimostri (usando la tecnica più opportuna) se da (1) e dalle seguenti osservazioni:
\begin{itemize}
    \item Ci sono studenti ansiosi.
    \item Ci sono studenti che hanno studiato.
\end{itemize}
sia corretto inferire che nessuno studente supererà alcun esame.
\subsection{Inferenza B}
Si dimostri (usando la tecnica più opportuna) se da (1) e dalle seguenti osservazioni:
\begin{itemize}
    \item Non tutti gli studenti sono ansiosi.
    \item Ogni studente ha studiato per almeno un esame.
\end{itemize}
sia corretto inferire che tutti gli studenti supereranno almeno un esame.
\section{Inferenze}
\subsection{Inferenza A}
KB: \newline
$Studente(X_1)\land Ansioso(X_1)$ \newline
$Studente(X_2) \land Corso(C_1) \land HaStudiato(X_2, C_1)$ \newline
$Studente(X)\land Corso(C) \land (Ansioso(X) \lor \lnot HaStudiato(X,C)) \rightarrow \lnot Supera(X, C)$ \newline
$\alpha$:\newline
$(Studente(X) \land Corso(C)) \rightarrow \lnot Supera(X, C)$ \newline
\subsubsection{CNF}
$KB_{CNF}$:\newline
1. \;\;$ Studente(X_1) $,\newline
2. \;\;$ Ansioso(X_1) $,\newline
3. \;\;$ Studente(X_2) $,\newline
4. \;\;$ Corso(C_1) $,\newline
5. \;\;$ HaStudiato(X_2, C_1) $,\newline
6. \;\;$ \lnot Studente(X) \lor \lnot Corso(C) \lor \lnot Ansioso(X) \lor \lnot Supera(X, C) $,\newline
7. \;\;$ \lnot Studente(X) \lor \lnot Corso(C) \lor HaStudiato(X, C) \lor \lnot Supera(X, C) $,\newline
$\alpha_{CNF}$:\newline
$\lnot Studente(Z) \lor \lnot Corso(E) \lor \lnot Supera(Z, E)$\newline
$\lnot \alpha_{CNF}$:\newline
8. \;\;$ Studente(X_3) $,\newline
9. \;\;$ Corso(C_2) $,\newline
10. \;\;$ Supera(X_3,C_2) $,\newline
\subsubsection{Risoluzione}
Si evita la proposizionalizzazione di $KB_{CNF} \land \lnot \alpha_{CNF}$ poiché (oltre a creare molte clausole) si evince che non
ci siano clausole che portino a una clausola vuota. Tutte le clausole atomiche sono skolemizzate, e non esiste
una coppia di costanti di Skolem $X_k,C_k$ per cui si ha $Studente(X_k) \land Corso(C_k) \land Ansioso(X_k) \land Supera(X_k, C_k)$ oppure
$Studente(X_k) \land Corso(C_k) \land \lnot HaStudiato(X_k, C_k) \land Supera(X_k, C_k)$.
Si ha quindi $KB \not \models \alpha$.

\subsection{Inferenza B}
KB: \newline
$Studente(X_1)\land \lnot Ansioso(X_1)$ \newline
$Studente(X) \rightarrow (Corso(C_1(X)) \land HaStudiato(X, C_1(X)))$ \newline
$Studente(Y)\land Corso(C) \land (Ansioso(Y) \lor \lnot HaStudiato(Y,C)) \rightarrow \lnot Supera(X, C)$ \newline
$\alpha$:\newline
$Studente(Z) \rightarrow ( Corso(C_2(Z)) \land Supera(Z, C_2(Z)))$ \newline
\subsubsection{CNF}
$KB_{CNF}$:\newline
1. \;\;$Studente(S_1)$,\newline
2. \;\; $\lnot Ansioso(S_1)$,\newline
3. \;\;$\lnot Studente(X) \lor Corso(C_1(X))$,\newline
4. \;\;$\lnot Studente(X) \lor HaStudiato(X, C_1(X))$,\newline
5. \;\;$\lnot Studente(Y) \lor \lnot Corso(C) \lor \lnot Ansioso(Y) \lor \lnot Supera(Y, C)$, \newline
6. \;\;$\lnot Studente(Y) \lor \lnot Corso(C) \lor HaStudiato(Y, C) \lor \lnot Supera(Y, C)$ \newline
$\lnot \alpha_{CNF}$: \newline
7. \;\;$Studente(Z)$,\newline
8. \;\;$\lnot Corso(C_2(Z)) \lor \lnot Supera(Z, C_2(Z))$,\newline
Otteniamo tutte le clausole possibili partendo da $ KB_{CNF} \land \lnot \alpha_{CNF} $: \newline
9. (1., 3.)\;\;$ Corso(C_1(S_1)) $\newline
10. (1., 4.)\;\; $ HaStudiato(S_1, C_1(S_1)) $ \newline
11. (1., 5.)\;\; $ \lnot Corso(C) \lor \lnot Ansioso(S_1) \lor \lnot Supera(S_1, C) $, \newline
12. (1., 6.)\;\; $ \lnot Corso(C) \lor HaStudiato(S_1, C) \lor \lnot Supera(S_1, C) $, \newline
13. (3., 5.)\;\; $ \lnot Studente(X) \lor \lnot Studente(Y) \lor \lnot Ansioso(Y) \lor \lnot Supera(Y, C_1(X)) $, \newline
14. (3., 6.)\;\; $ \lnot Studente(X) \lor \lnot Studente(Y) \lor HaStudiato(Y, C_1(X)) \lor \lnot Supera(Y, C_1(X))$, \newline
15. (3., 7.)\;\; $ Corso(C_1(Z)) $, \newline
16. (4., 7.)\;\; $ HaStudiato(Z, C_1(Z)) $, \newline
17. (5., 7.)\;\; $ \lnot Corso(C) \lor \lnot Ansioso(Z) \lor \lnot Supera(Z, C) $, \newline
18. (6., 7.)\;\; $ \lnot Corso(C) \lor HaStudiato(Z, C) \lor \lnot Supera(Z, C) $, \newline
19. (1. 13.)\;\; $ \lnot Studente(Y) \lor \lnot Ansioso(Y) \lor \lnot Supera(Y, C_1(S_1)) $, \newline
20. (1. 13.)\;\; $ \lnot Studente(X) \lor \lnot Ansioso(S_1) \lor \lnot Supera(S_1, C_1(X)) $, \newline
21. (1. 14.)\;\; $ \lnot Studente(Y) \lor HaStudiato(Y, C_1(S_1)) \lor \lnot Supera(Y, C_1(S_1)) $, \newline
22. (1. 14.)\;\; $ \lnot Studente(X) \lor HaStudiato(S_1, C_1(X)) \lor \lnot Supera(S_1, C_1(X)) $, \newline
23. (1. 19./20.)\;\; $ \lnot Ansioso(S_1) \lor \lnot Supera(S_1, C_1(S_1)) $, \newline
24. (1. 21./22.)\;\; $  HaStudiato(S_1, C_1(S_1)) \lor \lnot Supera(S_1, C_1(S_1)) $, \newline
25. (indir. 7. 13.)\;\; $ \lnot Ansioso(Z) \lor \lnot Supera(Z, C_1(Z)) $, \newline
26. (indir. 7. 14.)\;\; $  HaStudiato(Z, C_1(Z)) \lor \lnot Supera(Z, C_1(Z)) $ \newline
\newline
Dopodiché non si possono più ottenere altre clausole, si ha quindi $KB \not \models \alpha$.
\section{Prover9 / Mace4}
\subsection{Inferenza A}
\subsubsection{Input}
\begin{verbatim}
formulas(assumptions).

exists x (Studente(x) & Ansioso(x)).
exists x exists c (Studente(x) & Corso(c) & HaStudiato(x, c)).
all x all c ((Studente(x) & Corso(c) & (Ansioso(x) | -HaStudiato(x, c)))
         -> -Supera(x,c)).

end_of_list.

formulas(goals).

all x all c ((Studente(x) & Corso(c)) -> -Supera(x,c)).

end_of_list.
\end{verbatim}
\subsubsection{Risultati}
Dall'esecuzione dell'input precedente si ottiene
\begin{verbatim}
    SEARCH FAILED
\end{verbatim}
\subsection{Inferenza B}
\subsubsection{Input}
\begin{verbatim}
formulas(assumptions).

exists x (Studente(x) & -Ansioso(x)).
all x (Studente(x) -> (exists c (Corso(c) & HaStudiato(x,c)))).
all x all c ((Studente(x) & Corso(c) &(Ansioso(x) | -HaStudiato(x, c))) 
        -> -Supera(x, c)).

end_of_list.

formulas(goals).

all x (Studente(x) -> (exists c (Corso(c) & Supera(x, c)))).

end_of_list.
\end{verbatim}
\subsubsection{Risultati}
Dall'esecuzione dell'input precedente si ottiene
\begin{verbatim}
    SEARCH FAILED
\end{verbatim}
\end{document}